% This document was generated by an LLM.

\documentclass{essay}

\title{Rigorous Derivation of the Total Derivative}
\author{}
\date{}

\begin{document}

\maketitle

\section{Introduction}

This document provides a rigorous proof for the total derivative of a multivariable function $F(x, y(x))$ with respect to $x$. To ensure a logical flow, we first establish the single-variable Mean Value Theorem and use it to construct the formal definition of multivariable differentiability.

\section{Prerequisites}

\subsection{The Single-Variable Mean Value Theorem}
To analyze changes in multiple dimensions, we must first rely on a fundamental property of one-dimensional continuous functions.

\begin{lemma}[Mean Value Theorem]
If a function $f(t)$ is continuous on the closed interval $[a, b]$ and differentiable on the open interval $(a, b)$, then there exists at least one point $c \in (a, b)$ such that:
\[ f(b) - f(a) = f'(c)(b - a) \]
\end{lemma}

\subsection{The Increment Theorem (Multivariable Differentiability)}
In single-variable calculus, differentiability simply requires the limit of the difference quotient to exist. In multivariable calculus, differentiability requires that the function can be closely approximated by a tangent plane, with an error that vanishes rapidly. We derive this required condition using the Mean Value Theorem.

\begin{theorem}[The Increment of a Function]
Let $F(x,y)$ be a function with continuous partial derivatives $\frac{\partial F}{\partial x}$ and $\frac{\partial F}{\partial y}$ on an open region containing the point $(x,y)$. Let $\Delta x$ and $\Delta y$ be small increments. The total change in the function, $\Delta F = F(x+\Delta x, y+\Delta y) - F(x,y)$, can be expressed as:
\[ \Delta F = \frac{\partial F}{\partial x}\Delta x + \frac{\partial F}{\partial y}\Delta y + \varepsilon_1 \Delta x + \varepsilon_2 \Delta y \]
where $\varepsilon_1, \varepsilon_2 \to 0$ as $(\Delta x, \Delta y) \to (0,0)$.
\end{theorem}

\begin{proof}
Consider the total change $\Delta F = F(x+\Delta x, y+\Delta y) - F(x,y)$. We can break this diagonal movement into two orthogonal steps by introducing an intermediate point, $F(x+\Delta x, y)$:
\[ \Delta F = [F(x+\Delta x, y+\Delta y) - F(x+\Delta x, y)] + [F(x+\Delta x, y) - F(x,y)] \]

In the first bracket, only the $y$-variable changes. In the second, only the $x$-variable changes. We apply the single-variable Mean Value Theorem to each bracket:
\begin{enumerate}
    \item For the change in $x$, there exists a point $c_1$ between $x$ and $x+\Delta x$ such that:
    \[ F(x+\Delta x, y) - F(x,y) = \frac{\partial F}{\partial x}(c_1, y) \Delta x \]
    \item For the change in $y$, there exists a point $c_2$ between $y$ and $y+\Delta y$ such that:
    \[ F(x+\Delta x, y+\Delta y) - F(x+\Delta x, y) = \frac{\partial F}{\partial y}(x+\Delta x, c_2) \Delta y \]
\end{enumerate}

Substituting these back yields:
\[ \Delta F = \frac{\partial F}{\partial x}(c_1, y) \Delta x + \frac{\partial F}{\partial y}(x+\Delta x, c_2) \Delta y \]

We want this expression in terms of the derivatives evaluated at our starting point $(x,y)$. We define the error terms $\varepsilon_1$ and $\varepsilon_2$ as the difference between the derivative at the intermediate points and the derivative at the origin point $(x,y)$:
\[ \varepsilon_1 = \frac{\partial F}{\partial x}(c_1, y) - \frac{\partial F}{\partial x}(x, y) \implies \frac{\partial F}{\partial x}(c_1, y) = \frac{\partial F}{\partial x}(x, y) + \varepsilon_1 \]
\[ \varepsilon_2 = \frac{\partial F}{\partial y}(x+\Delta x, c_2) - \frac{\partial F}{\partial y}(x, y) \implies \frac{\partial F}{\partial y}(x+\Delta x, c_2) = \frac{\partial F}{\partial y}(x, y) + \varepsilon_2 \]

Substituting these definitions into our equation for $\Delta F$:
\[ \Delta F = \left[ \frac{\partial F}{\partial x}(x, y) + \varepsilon_1 \right] \Delta x + \left[ \frac{\partial F}{\partial y}(x, y) + \varepsilon_2 \right] \Delta y \]
\[ \Delta F = \frac{\partial F}{\partial x}\Delta x + \frac{\partial F}{\partial y}\Delta y + \varepsilon_1 \Delta x + \varepsilon_2 \Delta y \]

Because we assumed the partial derivatives are continuous, and because $c_1 \to x$ and $c_2 \to y$ as $\Delta x, \Delta y \to 0$, it must follow that $\varepsilon_1 \to 0$ and $\varepsilon_2 \to 0$. This concludes the proof of the Increment Theorem.
\end{proof}

\section{The Total Derivative Proof}

With the definition of multivariable differentiability established, we can now prove the chain rule for a path parameterized by $x$.

\begin{theorem}[Total Derivative]
Let $F(x,y)$ be continuously differentiable, and let $y = y(x)$ be a differentiable function of $x$. The total derivative of $F$ with respect to $x$ is given by:
\[ \frac{d}{dx}[F(x, y(x))] = \frac{\partial F}{\partial x} \frac{dx}{dx} + \frac{\partial F}{\partial y} \frac{dy}{dx} \]
\end{theorem}

\begin{proof}
Let $x$ undergo a small change $\Delta x$. Because $y$ is a function of $x$, this induces a corresponding change $\Delta y = y(x+\Delta x) - y(x)$.

From our prerequisites, because $F$ is continuously differentiable, the total change $\Delta F$ is exactly:
\[ \Delta F = \frac{\partial F}{\partial x}\Delta x + \frac{\partial F}{\partial y}\Delta y + \varepsilon_1 \Delta x + \varepsilon_2 \Delta y \]

Divide the entire equation by $\Delta x$ to evaluate the rate of change:
\[ \frac{\Delta F}{\Delta x} = \frac{\partial F}{\partial x} + \frac{\partial F}{\partial y}\frac{\Delta y}{\Delta x} + \varepsilon_1 + \varepsilon_2 \frac{\Delta y}{\Delta x} \]

Now, we take the limit as $\Delta x \to 0$:
\[ \lim_{\Delta x \to 0} \frac{\Delta F}{\Delta x} = \lim_{\Delta x \to 0} \left[ \frac{\partial F}{\partial x} + \frac{\partial F}{\partial y}\frac{\Delta y}{\Delta x} + \varepsilon_1 + \varepsilon_2 \frac{\Delta y}{\Delta x} \right] \]

We evaluate the limits of the individual components:
\begin{itemize}
    \item $\lim_{\Delta x \to 0} \frac{\Delta F}{\Delta x} = \frac{dF}{dx}$ by definition.
    \item $\lim_{\Delta x \to 0} \frac{\Delta y}{\Delta x} = \frac{dy}{dx}$ because $y(x)$ is differentiable.
    \item Because $y(x)$ is differentiable, it is continuous. Therefore, as $\Delta x \to 0$, $\Delta y \to 0$.
    \item Because $(\Delta x, \Delta y) \to (0,0)$, our Increment Theorem dictates that $\varepsilon_1 \to 0$ and $\varepsilon_2 \to 0$.
\end{itemize}

Substituting these limits into the equation:
\[ \frac{dF}{dx} = \frac{\partial F}{\partial x} + \frac{\partial F}{\partial y}\left(\frac{dy}{dx}\right) + 0 + 0 \cdot \left(\frac{dy}{dx}\right) \]
\[ \frac{dF}{dx} = \frac{\partial F}{\partial x} + \frac{\partial F}{\partial y}\frac{dy}{dx} \]

Noting that $\frac{dx}{dx} = 1$, we can express this in the requested symmetric form:
\[ \frac{dF}{dx} = \frac{\partial F}{\partial x}\frac{dx}{dx} + \frac{\partial F}{\partial y}\frac{dy}{dx} \]
\end{proof}

\newpage

\section{Self-Check Questions}

\textbf{Q: What specific condition on the partial derivatives of $F(x,y)$ guarantees that the error terms $\varepsilon_1$ and $\varepsilon_2$ vanish in the limit?} \\
\textbf{A:} The partial derivatives must be continuous (i.e., the function is of class $C^1$). This continuity ensures that the derivative evaluated at an intermediate point $(c_1, y)$ smoothly approaches the derivative at the exact point $(x,y)$ as $\Delta x \to 0$.

\vspace{1em}
\textbf{Q: During the derivation of the increment $\Delta F$, what is the mathematical purpose of introducing the intermediate point $F(x+\Delta x, y)$?} \\
\textbf{A:} It breaks the simultaneous diagonal shift into two distinct, orthogonal steps (one strictly parallel to the $x$-axis, one strictly parallel to the $y$-axis). This isolation is required to apply the single-variable Mean Value Theorem.

\vspace{1em}
\textbf{Q: When evaluating $\lim_{\Delta x \to 0}$, how do we formally justify that $\Delta y \to 0$ as well?} \\
\textbf{A:} The theorem requires $y = y(x)$ to be a differentiable function. Differentiability strictly implies continuity. By the definition of continuity, a vanishingly small change in the input ($\Delta x \to 0$) forces a vanishingly small change in the output ($\Delta y \to 0$).

\vspace{1em}
\textbf{Q: State the fundamental difference between the partial derivative $\frac{\partial F}{\partial x}$ and the total derivative $\frac{dF}{dx}$ in the context of this theorem.} \\
\textbf{A:} The partial derivative $\frac{\partial F}{\partial x}$ measures the rate of change of $F$ assuming all other variables (like $y$) are held strictly constant. The total derivative $\frac{dF}{dx}$ measures the complete rate of change of $F$ as $x$ varies, explicitly accounting for the fact that $y$ is forced to change alongside $x$ because $y = y(x)$.

\end{document}
